% =====================================================
% 题型：投篮独立重复试验与分布列及条件概率解答题 (q_stats_basketball_shooting.tex)
% 知识板块：概率分布与条件概率
% 主思维：结构归纳与套用
% 副思维：量词意识、卷面表达、充分必要
% 错因标签：分布列首末项计算错误、事件定义不清
% 讲义用途：概率过程表达
% =====================================================
% 难度：★★ (中档)
% =====================================================
\newcommand{\qStatsBasketballShootingStem}{%
  设整数 \(N \geqslant 2\)。某同学用一个球进行投篮练习，至多投篮 \(N\) 次，当且仅当投中 1 次时或 \(N\) 次均未投中时，停止练习。设该同学每次投中的概率为 \(p \ (0 < p < 1)\)，各次投中与否相互独立，记 \(X\) 为停止练习时该同学的投篮次数。\\
  (1) 当 \(N=4\)，\(p=\dfrac{1}{3}\) 时，求 \(X\) 的分布列；\\
  (2) 设 \(k\)，\(m\) 均为自然数。\\
  (i) 当 \(k \leqslant N-1\) 时，求 \(P(X > k)\)；\\
  (ii) 当 \(k+m \leqslant N-1\) 时，证明：\(P(X > k+m \mid X > k) = P(X > m)\)。%
}

\newcommand{\qStatsBasketballShootingAnswer}{%
  (1) 分布列见解析； (2)(i) \(P(X>k) = (1-p)^k\)； (ii) 见解析。%
}

\newcommand{\qStatsBasketballShootingSolution}{%
  【核心思维】\\
  第一步，明确停止练习的标志是“首次投中”或“投满 \(N\) 次都未中”，因此本质上是几何分布的截断；第二步，计算有限项的分步概率，列出分布列；第三步，将 \(X > k\) 的条件翻译为“前 \(k\) 次投篮均未命中”，求得其通项公式；第四步，利用条件概率定义化简比值，证明概率等价性（即几何分布的无记忆性在截断情况下的体现）。\\
  【标准作答与步骤分析】\\
  解：(1) 记每次投中为事件 \(M\)，每次未投中为事件 \(\overline{M}\)，且未投中概率 \(q = 1 - p = \dfrac{2}{3}\)。\\
  由题意，随机变量 \(X\) 的可能取值为 \(1, 2, 3, 4\)。\\
  当 \(X = 1\) 时，表示第 1 次投中：\\
    \(P(X=1) = p = \dfrac{1}{3}\)；\\
  当 \(X = 2\) 时，表示第 1 次未中，第 2 次投中：\\
    \(P(X=2) = qp = \dfrac{2}{3} \times \dfrac{1}{3} = \dfrac{2}{9}\)；\\
  当 \(X = 3\) 时，表示前 2 次未中，第 3 次投中：\\
    \(P(X=3) = q^2 p = \left(\dfrac{2}{3}\right)^2 \times \dfrac{1}{3} = \dfrac{4}{27}\)；\\
  当 \(X = 4\) 时，表示前 3 次均未投中，不管第 4 次投中与否，练习均停止：\\
    \(P(X=4) = q^3 = \left(\dfrac{2}{3}\right)^3 = \dfrac{8}{27}\)。\\
  所以，\(X\) 的概率分布列为：\\
  \[
  \begin{array}{c|cccc}
  X & 1 & 2 & 3 & 4 \\
  \hline
  P & \dfrac{1}{3} & \dfrac{2}{9} & \dfrac{4}{27} & \dfrac{8}{27} \\
  \end{array}
  \]
  (2) (i) 当 \(k \leqslant N-1\) 时，事件 \(X > k\) 发生，意味着停止练习的次数大于 \(k\)，\\
  即前 \(k\) 次投篮中，没有一次投中，因此前 \(k\) 次均未投中。\\
  由于各次投篮相互独立，所以其概率为：\\
  \(P(X > k) = (1 - p)^k\)。\\
  (ii) 证明：当 \(k+m \leqslant N-1\) 时，根据条件概率公式：\\
  \(P(X > k+m \mid X > k) = \dfrac{P(X > k+m \cap X > k)}{P(X > k)}\)。\\
  因为 \(k+m > k\)，所以事件 \(X > k+m\) 包含在 \(X > k\) 中，\\
  因此 \(P(X > k+m \cap X > k) = P(X > k+m)\)。\\
  由 (i) 中的公式代入得：\\
  \(P(X > k+m) = (1 - p)^{k+m}\)，\\
  \(P(X > k) = (1 - p)^k\)。\\
  所以：\\
  \(P(X > k+m \mid X > k) = \dfrac{(1 - p)^{k+m}}{(1 - p)^k} = (1 - p)^m\)。\\
  因为 \(m \leqslant k+m \leqslant N-1\)，由 (i) 亦有：\\
  \(P(X > m) = (1 - p)^m\)。\\
  故 \(P(X > k+m \mid X > k) = P(X > m)\) 成立。%
}

\newcommand{\qStatsBasketballShootingDiagnosis}{%
  第一问中，学生极易在 \(X=4\) 的概率计算上出错（常错写为前 3 次未中、第 4 次中或不中，而多乘了 \(p\) 或 \(q\) 的系数）。必须注意，在截断几何分布中，最后一项的概率包含投满未中的情况，因此直接等于前 \(N-1\) 次均不中的概率 \(q^{N-1}\)。%
}

\newcommand{\qStatsBasketballShootingTeachingNotes}{%
  本题考查概率的核心概念。第二问本质上是在截断几何分布中证明其“无记忆性”（Memoryless Property）。讲评时，可以对比几何分布的原生无记忆性，强调只要限制次数 \(k+m \leqslant N-1\)，此性质在截断后依然成立，借此加深对条件概率的物理理解。%
}
